% ======================================================================
% RUN 57 - P-M3: the clause-(i) difference theorem, PORTED BY RE-DERIVATION
% Target: item (M3) of appendix_e2a_omega.tex sec:an17-modulo, i.e. the
% difference bound eq:an17-clausei = clause (i) of hyp:hadamard-uniform
% (Option B ball form; a fortiori hyp:hadamard-uniform-omega clause (i)).
% Provenance: the banked LEMMA-II clause-(i) theorem (ISSUES.md l.4100
% region: C_H = C_sob^3 kappa, Gronwall/transport, exact Christoffel
% difference, r^{-2} kernel L^1_loc at n=4, state-independent). Per the
% run-48 law it is RE-DERIVED here in full, not transcribed; every
% constant is displayed. Paper-preamble-compatible fragment; compile
% check via port_m3_wrap.tex. ZERO writes to the three protected files.
% Labels are prefixed port: to avoid collisions.
% ======================================================================

\subsection*{P-M3. Clause (i), difference form: the dual-Lipschitz
singular-part theorem over the full $H^s$ ball, with constants}

\paragraph{Setting (the paper's own).} $n=4$, $\mathcal M\cong\mathbb
R_t\times\Sigma^3$ with the fixed $g_0$-atlas and compacts $K\subset
K'\subset K''$ of the appendix Setup; $|\cdot|$ and $\partial$ are
chart-Euclidean; $r:=|x-y|$; $t$ is the fixed $g_0$-time function and
$t_\Delta:=t(x)-t(y)$. Fix $s>n/2+6=8$ and the ball
\[
  U:=\bigl\{g\in H^s:\ \|g-g_0\|_{H^s(K'')}\le\varepsilon_0\bigr\},
\]
with $\varepsilon_0$ small enough that (u1) $\inf_U\min_{K''}|\det g|
=:d_U>0$; (u2) the $C^0$-image of $U$ has radius $<\varepsilon_{\mathrm
{GH}}$ (global hyperbolicity, \cite{ChruscielGrant2012} --- the
hypothesis' own smallness); (u3) the cone-transversality constants
$c_*,c_t$ of Lemma~\ref{port:lem-sigma}(iii) below are positive. This is
exactly the class $\mathcal U$ of Hyp.~\ref{hyp:hadamard-uniform} (Option B);
by Lem.~\ref{an17:lem-embed}, $U_\omega\subseteq U$ once $\varepsilon_a\le
\varepsilon_0/C_{\mathrm{emb}}$, so everything below holds verbatim over
$U_\omega$ with the same constants.

$H_g$ is the appendix's own singular part: the $N{=}2$ truncated glued
Hadamard parametrix of Lem.~\ref{an17:G},
\begin{equation}\label{port:eq-Hdef}
  H_g=\sum_{a\le N(K)}\theta_a\,Z^{(2),a}_g,\qquad
  Z^{(2),a}_g(x,y)=\frac{1}{8\pi^2}\Bigl[\frac{\Delta_g^{1/2}(x,y)}
  {\sigma_{g,\epsilon}(x,y)}
  +v^{(2)}_g(x,y)\,\log\frac{\sigma_{g,\epsilon}(x,y)}{\ell^2}\Bigr],
\end{equation}
$v^{(2)}_g=\sum_{k=0}^{2}v_k[g]\,\sigma_g^k$, on the bi-collar
$V_{\delta_c}=\{x\in K,\ r\le\delta_c\}$ of the ball-uniform
convex-normal cover (radius $\delta_c$: Chen--LeFloch /
Cheeger--Gromov--Taylor, cited at its printed site Lem.~\ref{an17:G}, THEOREM;
the cutoffs $\theta_a$ are $g$-\emph{independent}, built from the
$g_0$-atlas partition of unity, $\|\theta_a\|_{C^2}\le c_\theta$,
overlap multiplicity $m(K)$). Here $\sigma_g$ is the world function,
$\Delta_g^{1/2}$ the van Vleck determinant, $v_k[g]$ the Hadamard
coefficients of the transport recursion (\cite[Thm.~2.1]{Moretti2003}),
$\ell$ the single global length scale (Lem.~\ref{an17:G}(ii)), and
$\sigma_{g,\epsilon}=\sigma_g+2i\epsilon t_\Delta+\epsilon^2$ the
prescription of \cite{KayWald1991}, $\epsilon\downarrow0$. The pairing
$H_g(u\otimes v)$ is the $\epsilon\downarrow0$ limit, which exists per
metric on the cited footing (\cite{Moretti2003,KayWald1991}; this is the
single-metric half, established as Lem.~\ref{lem:an17-clausei-single}); every bound
below is uniform in $\epsilon\in(0,1]$, so it passes to the limit.

\begin{theorem}[Clause (i), difference form, over the full $H^s$ ball;
every constant displayed]\label{port:thm-M3}
Let $n=4$, $s>n/2+6$, $U$ as above, and set
$\alpha:=\min\bigl(\tfrac12,\tfrac{s-8}{2}\bigr)\in(0,\tfrac12]$. There
is a constant
\[
  C_H(K)\;=\;C_{\mathrm{sob}}(K)^3\,\kappa(g_0,K),
\]
with $C_{\mathrm{sob}}(K)$ and $\kappa(g_0,K)$ the displayed constants
\eqref{port:eq-csob} and \eqref{port:eq-kappa} below, such that for all
$g,g'\in U$ and all $u,v$ supported in $K$,
\begin{equation}\label{port:eq-clausei}
  \bigl|(H_g-H_{g'})(u\otimes v)\bigr|
  \;\le\;C_H(K)\,\|g-g'\|_{H^s}\,\|u\|_{H^{s-2}}\,\|v\|_{H^{s-2}}.
\end{equation}
The bound is ball-uniform (it depends on $U$ only through
$(g_0,\varepsilon_0,s,K,\delta_c,\ell)$), state-independent (no state
enters \eqref{port:eq-Hdef}), and holds a fortiori over
$U_\omega\times U_\omega$. Sharper form proved en route: one test slot
needs only $\sup$-norm,
$\;|(H_g-H_{g'})(u\otimes v)|\le C_{\mathrm{sob}}\,\kappa\,
\|g-g'\|_{H^s}\|u\|_{C^0}\|v\|_{C^{1,\alpha}}$
(one embedding, on the metric difference alone).
\end{theorem}

\paragraph{The constant inventory.} All constants are finite and
ball-uniform; each is defined at its first use and collected here.
\begin{equation}\label{port:eq-csob}
  C_{\mathrm{sob}}(K):=\max\bigl\{
  C^{\mathrm{Mor}}_{s\to(6,\alpha)}(K''),\;
  C^{\mathrm{Mor}}_{(s-2)\to(1,\alpha)}(K)\bigr\},
\end{equation}
the two Morrey--Sobolev embedding constants
$\|\cdot\|_{C^{6,\alpha}(K'')}\le C^{\mathrm{Mor}}_{s\to(6,\alpha)}
\|\cdot\|_{H^{s}(K'')}$ (legal: $s\ge8+2\alpha>8$, $n=4$) and
$\|\cdot\|_{C^{1,\alpha}(K)}\le C^{\mathrm{Mor}}_{(s-2)\to(1,\alpha)}
\|\cdot\|_{H^{s-2}(K)}$ (legal: $s-2>6>3+\alpha$); Adams--Fournier.
The three factors of $C_{\mathrm{sob}}^3$ are the three norm
conversions of \eqref{port:eq-clausei}: one on the metric difference
($\|h\|_{C^6}\le C_{\mathrm{sob}}\|h\|_{H^s}$, $h:=g-g'$), one on each
test slot. Ball-uniform geometry constants:
$\Lambda_U:=1+\sup_{g\in U}\|g\|_{C^{6,\alpha}(K'')}
\le1+C_{\mathrm{sob}}(\|g_0\|_{H^s(K'')}+\varepsilon_0)$;
$d_U$ and $\delta_c$ as above; $|K|$, $|K''|$ chart volumes;
$\Lambda_{\mathrm{vol}}:=\sup_U\sup_{K''}\sqrt{|\det g|}\le
c\,\Lambda_U^{2}$. The masses of the two singular profiles
($n=4$; the $r^{-2}$ kernel is where $L^1_{\mathrm{loc}}$ at $n=4$
enters):
\begin{equation}\label{port:eq-masses}
  M_2:=\sup_{x\in K}\int_{r\le\delta_c}r^{-2}\,dy
  =\tfrac{|S^3|}{2}\delta_c^{2}=\pi^2\delta_c^2,\qquad
  M_{\log}:=\sup_{x\in K}\int_{r\le\delta_c}\Bigl(1+\bigl|\log
  \tfrac{\Lambda_\sigma r^2}{\ell^2}\bigr|\Bigr)dy<\infty .
\end{equation}

\begin{lemma}[L1: Christoffel difference --- exact identity, then the
constant]\label{port:lem-gamma}
For $g,g'\in U$, $h:=g-g'$, with $\nabla'$ the $g'$-connection,
\begin{equation}\label{port:eq-gammaid}
  \delta\Gamma^{\lambda}{}_{\mu\nu}
  :=\Gamma[g]^{\lambda}{}_{\mu\nu}-\Gamma[g']^{\lambda}{}_{\mu\nu}
  =\tfrac12\,g^{\lambda\rho}\bigl(\nabla'_{\mu}h_{\rho\nu}
  +\nabla'_{\nu}h_{\rho\mu}-\nabla'_{\rho}h_{\mu\nu}\bigr),
\end{equation}
an exact tensor identity (no remainder). Consequently, for $0\le k\le5$,
\begin{equation}\label{port:eq-gammaconst}
  \|\delta\Gamma\|_{C^k(K'')}\le\kappa_\Gamma(k)\,\|h\|_{C^{k+1}(K'')},
  \qquad
  \kappa_\Gamma(k):=c_k(n)\,\Lambda_U^{\,2k+7}\,d_U^{-2(k+1)} .
\end{equation}
\end{lemma}
\begin{proof}
The difference of two torsion-free metric connections is a tensor;
evaluating $\nabla' g=\nabla' (g'+h)=\nabla' h$ in the Koszul formula
for $\Gamma[g]$ written in $\nabla'$-form gives \eqref{port:eq-gammaid}
directly (or: verify in $g'$-normal coordinates at a point, where
$\Gamma[g']=0$ and \eqref{port:eq-gammaid} is the Koszul formula for
$h$; both sides are tensors, so the identity is global). For
\eqref{port:eq-gammaconst}: $g^{-1}=\mathrm{cof}(g)^{\mathsf T}/\det g$
with $\mathrm{cof}$ polynomial of degree $n-1=3$ in the entries of
$g$, so $\|g^{-1}\|_{C^k(K'')}\le c_k\Lambda_U^{3+k}d_U^{-(k+1)}$
(Leibniz on the quotient, $|\det g|\ge d_U$); $\nabla'h=\partial h+
\Gamma'\!\cdot\!h$ with $\|\Gamma'\|_{C^k}\le c\Lambda_U^{k+4}
d_U^{-(k+1)}$ by the same count at the single metric $g'$. Multiply
out and absorb the numerical factors into $c_k(n)$.
\end{proof}

\begin{lemma}[L2: world-function difference at the Gronwall rate; the
rescaled normal form; uniform cone transversality]\label{port:lem-sigma}
There are ball-uniform constants $\Lambda_\sigma,\kappa_\sigma,
\Lambda_S,\kappa_S,c_*,c_t>0$, displayed in the proof, such that on
$V_{\delta_c}$ (per convex-normal chart $a$, all statements uniform in
$a$):
\begin{enumerate}[label=\textup{(\roman*)},leftmargin=*,nosep]
\item (rescaled normal form) $\sigma_g(x,x+r\omega)=r^2S_g(x,\omega,r)$,
  $|\omega|=1$, with $S_g\in C^3$,
  $\|S_g\|_{C^3(K\times S^3\times[0,\delta_c])}\le\Lambda_S$, and
  $S_g(x,\omega,0)=\tfrac12\,q_{g,x}(\omega,\omega)$ where
  $q_{g,x}$ is the metric quadratic form at $x$; in particular
  $|\partial^j_{(x,y)}\sigma_g|\le\Lambda_\sigma r^{(2-j)_+}$, $j\le2$;
\item (difference legs, Lipschitz, real side) for $j=0,1,2$,
  \[
  \bigl|\partial^j_{(x,y)}(\sigma_g-\sigma_{g'})\bigr|
  \le\kappa_\sigma\,\|h\|_{C^{j+1}(K'')}\,r^{2-j},
  \qquad\text{equivalently}\quad
  \|S_g-S_{g'}\|_{C^j}\le\kappa_S\,\|h\|_{C^{j+1}(K'')};
  \]
\item (uniform transversality) for every $\tau\in[0,1]$ the interpolant
  $S_\tau:=(1-\tau)S_{g'}+\tau S_g$ satisfies, on the cone collar
  $\{|S_\tau|\le\theta_*\}$, $|\nabla_\omega S_\tau|\ge c_*>0$ and
  $|t_\Delta|\ge c_t\,r$ with $\mathrm{sign}\,t_\Delta$ constant on
  each (future/past) sheet; $\theta_*=\theta_*(d_U,\Lambda_U)>0$.
\end{enumerate}
\end{lemma}
\begin{proof}
(i) In the chart, the $g$-exponential flow $\Phi^t_g(x,\xi)$ solves
$\ddot\gamma=-\Gamma[g](\dot\gamma,\dot\gamma)$; on the compact shell
$\{x\in K,|\xi|\le2\delta_c\}$ its $C^3$-norm in $(x,\xi,t)_{t\in[0,1]}$
is bounded by the Picard/Gronwall constant
$\kappa_{\mathrm{fl}}:=c\,\exp\bigl(c\,\delta_c\,
\|\Gamma\|_{C^3(K'')}\bigr)\le
c\exp\bigl(c\,\delta_c\,\Lambda_U^{7}d_U^{-4}\bigr)$
(coefficients $\Gamma[g]\in C^{4,\alpha}$ since $g\in C^{6,\alpha}$;
three derivatives of the flow consume $\|\Gamma\|_{C^3}$). With
$\xi=\exp_x^{-1}(y)$ (well defined and $C^3$, Jacobian
$\ge\tfrac12$ on $V_{\delta_c}$ after the one ball-uniform shrink of
$\delta_c$ built into Lem.~\ref{an17:G}) and
$\sigma_g(x,y)=\tfrac12 q_{g,x}(\xi,\xi)$, homogeneity of the geodesic
flow in $(\xi,t)$ gives $\xi=r\,\Xi_g(x,\omega,r)$ with
$\Xi_g\in C^3$, $\Xi_g(x,\omega,0)=\omega$; set
$S_g:=\tfrac12 q_{g,x}(\Xi_g,\Xi_g)$. Then $\Lambda_S\le
c\,\Lambda_U\kappa_{\mathrm{fl}}^2$ and
$\Lambda_\sigma\le c\,\Lambda_S$.
(ii) The flows of $g,g'$ differ by a linear-inhomogeneous
(variational) ODE with source $\delta\Gamma(\dot\gamma,\dot\gamma)$:
Gronwall gives
$\|\Phi_g-\Phi_{g'}\|_{C^j(\text{shell})}\le
\kappa_{\mathrm{fl}}^2\,\|\delta\Gamma\|_{C^j}$, $j\le2$, hence
$\|\Xi_g-\Xi_{g'}\|_{C^j}\le c\,\kappa_{\mathrm{fl}}^3
\|\delta\Gamma\|_{C^j}$ (inverse-function difference at Jacobian
$\ge\tfrac12$), and with \eqref{port:eq-gammaconst},
$\kappa_S:=c\,\Lambda_U\kappa_{\mathrm{fl}}^{3}
\bigl(1+\kappa_\Gamma(3)\bigr)$,
$\kappa_\sigma:=c\,\kappa_S$. The $r$-grading in the unscaled form is
the chain rule applied to $r^2S(x,\omega,r)$, $\omega=(y-x)/r$: each
$(x,y)$-derivative costs one factor $r$ at worst.
(iii) At $r=0$: $\nabla_\omega S_\tau=q_{\tau,x}(\omega,\cdot)$ with
$q_\tau:=(1-\tau)q_{g'}+\tau q_g$ --- again a Lorentzian metric in the
convex ball $U$, $|\det q_\tau|\ge d_U'>0$ for $\varepsilon_0$ small
(continuity of $\det$ on the $C^0$-ball). On the set
$\{S_\tau=0,|\omega|=1\}$, $\omega$ is $q_\tau$-null, and for a
Lorentzian form with $|\det|\ge d_U'$, $\|q\|\le\Lambda_U$ one has
$|q(\omega,\cdot)|\ge c(d_U',\Lambda_U)>0$ on unit null vectors
(else $\omega$ would be a null eigenvector of $q$ with vanishing
eigenvalue, contradicting $|\det q|\ge d_U'$); perturbing in
$r\le\delta_c$ and in $|S_\tau|\le\theta_*$ costs a factor
$\tfrac12$ for $\theta_*,\delta_c$ small ball-uniformly:
$c_*:=\tfrac12c(d_U',\Lambda_U)$. For $c_t$: on the $q_\tau$-null
cone, $|\omega^0|\ge c(d_U',\Lambda_U)>0$ (time component of unit
null vectors is bounded below when the cones are pinched between two
$g_0$-cones, which is (u2)-smallness), and
$t_\Delta=-r\,\omega^0\,\partial_0t+O(\Lambda_S r^2)$, so
$c_t:=\tfrac12c(d_U',\Lambda_U)\inf_{K''}|\partial_0 t|$ after one
more shrink of $\delta_c$; the sign is constant per sheet because
$\omega^0$ has a fixed sign on each component of the null cone. The
time leg extends from the collar to the full timelike side
$\{S_\tau\le\theta_*\}$ by the same linear algebra: on unit
vectors with $q_\tau(\omega,\omega)\le\theta_*$, a time component
$|\omega^0|<c'$ would force $q_\tau(\omega,\omega)\ge
c(d_U',\Lambda_U)-O(c')>\theta_*$ (the spatial block is positive
definite with determinant bounded below), a contradiction; so
$|\omega^0|\ge c'(d_U',\Lambda_U)>0$, hence $|t_\Delta|\ge c_t r$
with per-sheet fixed sign, on all of $\{S_\tau\le\theta_*\}$,
$r\le\delta_c$.
\end{proof}

\begin{lemma}[L3: van Vleck and Hadamard-coefficient legs ---
transport/Gronwall, sup and difference]\label{port:lem-vk}
With $\Lambda_S,\kappa_S$ as in L2, there are displayed ball-uniform
$\Lambda_\Delta,\kappa_\Delta,\Lambda_v,\kappa_v$ with, on
$V_{\delta_c}$:
\begin{enumerate}[label=\textup{(\roman*)},leftmargin=*,nosep]
\item $\|\Delta^{\pm1/2}_g\|_{C^{1,\alpha}(V)}\le\Lambda_\Delta$ and
  $\|\Delta^{1/2}_g-\Delta^{1/2}_{g'}\|_{C^{0,\alpha}(V)}
  \le\kappa_\Delta\|h\|_{C^{4}(K'')}$;
\item for $k\le2$: $\|v_k[g]\|_{C^{0,\alpha}(V)}\le\Lambda_v$ and
  $\|v_k[g]-v_k[g']\|_{C^{0}(V)}\le\kappa_v\|h\|_{C^{2k+2}(K'')}
  \le\kappa_v\|h\|_{C^{6}(K'')}$.
\end{enumerate}
\end{lemma}
\begin{proof}
(i) $\Delta_g(x,y)=-\det\bigl(-\partial_x\partial_y\sigma_g\bigr)
\bigl(|\det g(x)||\det g(y)|\bigr)^{-1/2}$; by L2(i)--(ii) the entries
$\partial_x\partial_y\sigma_g$ are $C^{1,\alpha}$ with norm
$\le c\Lambda_S$ and difference $\le c\kappa_S\|h\|_{C^{4}}$ in
$C^{0,\alpha}$ (one derivative above the $j=2$ leg, obtained by running
L2(ii) at $j=2$ on the $C^{1,\alpha}$ flow bounds --- the flow is $C^3$,
one order to spare); the determinant is a cubic polynomial, and on
$V_{\delta_c}$ the matrix is invertible with determinant
$\ge\tfrac12 d_U$ (coincidence value $\det(-\partial^2\sigma)=
|\det g|$, plus the ball-uniform shrink of $\delta_c$), so the square
roots are Lipschitz functions of the entries:
$\Lambda_\Delta:=c\,\Lambda_S^{4}d_U^{-3/2}$,
$\kappa_\Delta:=c\,\Lambda_S^{3}d_U^{-3/2}\kappa_S$.
(ii) The recursion of \cite[Thm.~2.1]{Moretti2003},
$2\nabla\sigma\!\cdot\!\nabla v_k+(\square\sigma-2n+4k)v_k
=-\square v_{k-1}$ ($v_{-1}$-convention anchoring
$v_0$ at $\Delta^{1/2}$-data), is for each $k$ a linear transport ODE
along the L2 geodesics: in the rescaled variable it reads
$\bigl(2r\partial_r+A_k(x,\omega,r)\bigr)v_k=f_k$ with
$A_k=\square\sigma-2n+4k+O(r)$, $|A_k|\le c\Lambda_S$, and source
$f_k$ a universal polynomial in the jets
$(\partial^{\le2}\sigma,\Delta^{\pm1/2},\partial^{\le2k+2}g,g^{-1})$.
Duhamel along $[0,r]$ with the regular-singular weight
$r^{-A_k/2}$-type factor (the indicial root is positive for
$k\ge0$, $n=4$: $\square\sigma\to2n$ at coincidence, so
$A_k\to4k>0$ for $k\ge1$ and the $k=0$ equation fixes
$v_0=\Delta^{1/2}$-transported data) gives the sup bound by Gronwall,
$\Lambda_v:=c\,e^{c\Lambda_S\delta_c}\,
P\bigl(\Lambda_S,\Lambda_\Delta,\Lambda_U,d_U^{-1}\bigr)$
($P$ a displayed universal polynomial; the $C^{0,\alpha}$ grade rides
the same bound one order down, legal since $g\in C^{6,\alpha}$ and
$f_2$ consumes exactly $\partial^{6}g\in C^{0,\alpha}$). The
difference leg is the same Duhamel applied to
$(v_k[g]-v_k[g'])$, whose ODE has coefficient difference
$\le c(\kappa_S+\kappa_\Gamma(1))\|h\|_{C^{2}}$ and source difference
$\le cP\cdot\|h\|_{C^{2k+2}}$ (linearity of $f_k$-assembly in its
top jet; lower-order factors estimated by the sup bounds already
obtained --- this is the linear-in-top-jet structure the appendix's
(N-a) block records):
$\kappa_v:=c\,e^{c\Lambda_S\delta_c}P\cdot
\bigl(1+\kappa_S+\kappa_\Gamma(1)\bigr)$. No step differentiates in
$g$; every difference is pair-Lipschitz on the real side.
\end{proof}

\begin{lemma}[L4: uniform cone pairing --- the analytic core]
\label{port:lem-pair}
Fix $x\in K$ and a phase $\rho=r^2S$ with $S\in\{S_\tau\}_{\tau\in[0,1]}$
as in L2 (so $\|S\|_{C^{3}}\le\Lambda_S$ --- L2(i) proves the $C^{3}$
flow bound --- transversality $(c_*,\theta_*)$,
time leg $(c_t,\Lambda_T:=\|t_\Delta/r\|_{C^1})$), and
$\rho_\epsilon:=\rho+2i\epsilon t_\Delta+\epsilon^2$. Let
$F(y)=r^{a}\Phi(\omega,r)$ in $y=x+r\omega$, $\operatorname{supp}F
\subset\{r\le\delta_c\}$. Then, uniformly in $\epsilon\in(0,1]$ and in
$\tau$:
\begin{enumerate}[label=\textup{(P\arabic*)},leftmargin=*,nosep]
\item if $a\ge0$ and $\Phi\in C^{0,\alpha}$:
  $\displaystyle\Bigl|\int\frac{F}{\rho_\epsilon}\,dy\Bigr|
  \le\kappa_{P}\,\|\Phi\|_{C^{0,\alpha}}\,\frac{\delta_c^{a+2}}{a+2}$;
\item if $a\ge2$ and $\Phi\in C^{1,\alpha}$:
  $\displaystyle\Bigl|\int\frac{F}{\rho_\epsilon^{2}}\,dy\Bigr|
  \le\kappa_{P}\,\|\Phi\|_{C^{1,\alpha}}\,\frac{\delta_c^{a}}{a}$;
\item if $a\ge0$ and $\Phi\in C^{0}$:
  $\displaystyle\int|F|\,\bigl|\log(\rho_\epsilon/\ell^2)\bigr|\,dy
  \le\kappa_{P}\,\|\Phi\|_{C^{0}}\,\delta_c^{a+2}
  \bigl(1+|\log\delta_c^2/\ell^2|\bigr)$;
\end{enumerate}
with the single displayed constant
\begin{equation}\label{port:eq-kpair}
  \kappa_{P}:=c_n\bigl(1+\Lambda_S+c_*^{-1}\bigr)^{4}
  \Bigl[\theta_*^{-2}\varkappa^{-2}+(\pi+4)+\tfrac{2}{\alpha}\theta_*^{\alpha}
  +\bigl(1+c_t^{-2}\bigr)\bigl(1+\Lambda_T\bigr)\Bigr],\qquad
  \varkappa:=\min\bigl(\tfrac12,\,\sqrt2\,c_t/\sqrt{\Lambda_S}\bigr).
\end{equation}
\end{lemma}
\begin{proof}
Rescale: $dy=r^3\,dr\,d\omega$, so
$\int F\rho_\epsilon^{-2}dy=\int_0^{\delta_c}r^{a-1}
\bigl[\int_{S^3}\Phi\,(S+i\eta)^{-2}d\omega\bigr]dr$ with
$\eta(\omega):=(2\epsilon t_\Delta+\epsilon^2)/r^2
=\eta_0T(\omega,r)+\epsilon^2/r^2$, $\eta_0:=2\epsilon/r$,
$T:=t_\Delta/r$ ($|T|\ge c_t$ with fixed sign per sheet on the cone
collar, L2(iii)); similarly $\int F\rho_\epsilon^{-1}dy
=\int_0^{\delta_c}r^{a+1}[\,\cdot\,]dr$ with $(S+i\eta)^{-1}$. The real
shift $\epsilon^2/r^2$: where $\epsilon^2\ge\theta_*r^2/2$, i.e.\
$r\le\epsilon\sqrt{2/\theta_*}$, split by the size of $r^{2}|S|$.
If $r^{2}|S|\le\epsilon^{2}/2$ then
$|\rho_\epsilon|\ge\operatorname{Re}\rho_\epsilon
=r^{2}S+\epsilon^{2}\ge\epsilon^{2}/2$; otherwise
$r^{2}\ge\epsilon^{2}/(2\Lambda_S)$ (since $|S|\le\Lambda_S$) and
the time leg --- valid on the whole set $\{S_\tau\le\theta_*\}$,
collar plus timelike interior, by L2(iii)'s extension --- gives
$|\rho_\epsilon|\ge|\operatorname{Im}\rho_\epsilon|
=2\epsilon|t_\Delta|\ge2\epsilon c_t r
\ge\sqrt2\,c_t\,\epsilon^{2}/\sqrt{\Lambda_S}$. Hence
$|\rho_\epsilon|\ge\epsilon^{2}\varkappa$ with
$\varkappa:=\min\bigl(\tfrac12,\,\sqrt2\,c_t/\sqrt{\Lambda_S}\bigr)$,
and $|\rho_\epsilon|\ge\varkappa\,\theta_*r^{2}/2$ on this range,
so the slice integrand is $\le\|\Phi\|_\infty
(2/\theta_*)^2\varkappa^{-2}|S^3|$ pointwise and this radial range
contributes $\le c\|\Phi\|_\infty\theta_*^{-2}\varkappa^{-2}
\delta_c^{a}/a$; elsewhere absorb
$\epsilon^2/r^2\le\theta_*/2$ into $S$ (transversality survives with
$c_*/2$, $\theta_*/2$). It remains to bound the slice integrals
$J_p:=\int_{S^3}\Phi(S+i\eta_0T)^{-p}d\omega$, $p=1,2$, uniformly in
$(r,\eta_0)$. Split $S^3$: on $\{|S|>\theta_*\}$,
$|J_p^{\mathrm{far}}|\le\|\Phi\|_\infty|S^3|\theta_*^{-p}$. On the cone
collar, the level flow of $S/|\nabla S|^2$ gives coordinates
$(z,\zeta)$, $z\in Z=\{S=0\}$, $S=\zeta$, $|\zeta|\le\theta_*$, with
Jacobian $J_z$, $\|J_z\|_{C^{1,\alpha}}\le c(1+\Lambda_S+c_*^{-1})^{4}$
(one further implicit-function differentiation, licensed by
$\|S\|_{C^{3}}\le\Lambda_S$ at $|\nabla S|\ge c_*$ --- the $p=2$
fiber Taylor split consumes $J_z\in C^{1,\alpha}$; the null sphere $Z$ has at most two sheets, the future and
past cones of the Lorentzian normal form, L2(iii)). Per fiber, with
$G:=\Phi J_z$ and $T$ frozen at its foot value $T_0$ ($|T_0|\ge c_t$,
sign fixed), $a:=\eta_0T_0\ne0$: the exact complex antiderivatives give the four
moment bounds, each uniform in $a\ne0$ ($L:=\theta_*$),
\begin{equation}\label{port:eq-moments}
  \Bigl|\int_{-L}^{L}\frac{d\zeta}{(\zeta+ia)^{2}}\Bigr|
  =\Bigl|\frac{-2L}{(L+ia)(-L+ia)}\Bigr|\le\frac{2}{L},\quad
  \Bigl|\int_{-L}^{L}\frac{\zeta\,d\zeta}{(\zeta+ia)^{2}}\Bigr|
  =\Bigl|\log\frac{L+ia}{-L+ia}+\frac{ia}{\zeta+ia}\Big|_{-L}^{L}\Bigr|
  \le\pi+2,
\end{equation}
\begin{equation}\label{port:eq-moment1}
  \Bigl|\int_{-L}^{L}\frac{d\zeta}{\zeta+ia}\Bigr|
  =\Bigl|\log\frac{L+ia}{-L+ia}\Bigr|\le\pi,\qquad
  \int_{-L}^{L}\frac{|\zeta|^{\beta}\,d\zeta}{|\zeta+ia|^{2}}
  \le\int_{-L}^{L}|\zeta|^{\beta-2}d\zeta
  =\tfrac{2}{\beta-1}L^{\beta-1}\ \ (\beta>1),
\end{equation}
all uniform in $a\ne0$ (the middle identity uses
$|L+ia|=|-L+ia|$, so the log has modulus $\le\pi$, its imaginary
part). For $p=2$ Taylor-split $G=G_0+\zeta G_0'+R$,
$|R|\le[G']_{\alpha}|\zeta|^{1+\alpha}$, and apply
\eqref{port:eq-moments}--\eqref{port:eq-moment1} with $\beta=1+\alpha$:
the fiber integral is $\le\bigl(\tfrac{2}{\theta_*}+\pi+2
+\tfrac{2}{\alpha}\theta_*^{\alpha}\bigr)\|G\|_{C^{1,\alpha}}$. For
$p=1$ split $G=G_0+(G-G_0)$, $|G-G_0|\le[G]_\alpha|\zeta|^\alpha$:
fiber $\le(\pi+\tfrac{2}{\alpha}\theta_*^{\alpha})\|G\|_{C^{0,\alpha}}$.
Unfreezing $T$: $|\eta(\zeta)-\eta_0T_0|\le\eta_0\Lambda_T|\zeta|/c_*$
and $|\zeta+i\eta|\ge\max(|\zeta|,\eta_0c_t)/2$ on the fiber, so the
freezing error for $p=2$ is
$\le c\|G\|_\infty\Lambda_T c_*^{-1}\int\eta_0|\zeta|
\max(|\zeta|,\eta_0c_t)^{-3}d\zeta\le
c\|G\|_\infty\Lambda_Tc_*^{-1}(1+c_t^{-2})$, and for $p=1$
$\le c\|G\|_\infty\Lambda_Tc_*^{-1}(1+c_t^{-1})$ --- uniform in
$\eta_0>0$. Integrate over $z\in Z$ ($|Z|\le c_n$) and the two sheets,
then radially: $\int_0^{\delta_c}r^{a-1}dr=\delta_c^{a}/a$ for (P2)
($a\ge2$; in unrescaled terms the integrand is dominated by
$r^{a}r^{-4}\cdot r^{3}=r^{a-1}$, the $n=4$ $L^1_{\mathrm{loc}}$ of the
$r^{-2}$ kernel at $a=2$, mass $M_2$ of \eqref{port:eq-masses}), and
$\int_0^{\delta_c}r^{a+1}dr=\delta_c^{a+2}/(a+2)$ for (P1). For (P3):
$|\log(\rho_\epsilon/\ell^2)|\le2|\log r|+|\log|S+i\eta||+
|{\log\ell^2}|+2\pi$ and $\int_{S^3}|\log|S+i\eta||\,d\omega\le
c(1+|\log\theta_*|)$ uniformly in $\eta$ (on fibers
$|\log|\zeta+ia||\le|\log|\zeta||+c$ for $|\zeta|\le\theta_*\le c$),
so the radial integral converges to the stated mass. Collect
\eqref{port:eq-kpair}.
\end{proof}

\begin{proof}[Proof of Theorem~\ref{port:thm-M3} (assembly; four legs)]
Fix $\epsilon\in(0,1]$ and a chart $a$; drop the chart index. With
$w:=\sigma_g-\sigma_{g'}$ and the linear interpolant
$\sigma_{\tau,\epsilon}:=(1-\tau)\sigma_{g',\epsilon}
+\tau\sigma_{g,\epsilon}=r^2S_\tau+2i\epsilon t_\Delta+\epsilon^2$
(same $t_\Delta$: the prescription is $g$-independent), the exact
calculus identities
$\;\sigma_{g,\epsilon}^{-1}-\sigma_{g',\epsilon}^{-1}
=-\int_0^1 w\,\sigma_{\tau,\epsilon}^{-2}\,d\tau\;$ and
$\;\log\sigma_{g,\epsilon}-\log\sigma_{g',\epsilon}
=\int_0^1 w\,\sigma_{\tau,\epsilon}^{-1}\,d\tau\;$ give
\[
  8\pi^2\bigl(Z^{(2)}_g-Z^{(2)}_{g'}\bigr)
  =\underbrace{\frac{\Delta^{1/2}_g-\Delta^{1/2}_{g'}}
  {\sigma_{g,\epsilon}}}_{\mathrm{(A)}}
  -\underbrace{\Delta^{1/2}_{g'}\!\int_0^1\!
  \frac{w\,d\tau}{\sigma_{\tau,\epsilon}^{2}}}_{\mathrm{(B)}}
  +\underbrace{\bigl(v^{(2)}_g-v^{(2)}_{g'}\bigr)
  \log\frac{\sigma_{g,\epsilon}}{\ell^2}}_{\mathrm{(C)}}
  +\underbrace{v^{(2)}_{g'}\!\int_0^1\!
  \frac{w\,d\tau}{\sigma_{\tau,\epsilon}}}_{\mathrm{(D)}} .
\]
Pair each leg with $\theta_a\,u(x)v(y)$, integrate in $y$ at fixed $x$
by L4, then in $x$. The graded weights, from L2--L3 (all differences
real-side, none through any collar):
(A) $F=(\Delta^{1/2}_g-\Delta^{1/2}_{g'})\theta_a u(x)v$: $a=0$,
$\|\Phi\|_{C^{0,\alpha}}\le c_\theta\kappa_\Delta\|h\|_{C^4}
\,|u(x)|\,\|v\|_{C^{0,\alpha}}$; apply (P1).
(B) $F=\Delta^{1/2}_{g'}w\,\theta_a u(x)v$: $w=r^2W$ with
$\|W\|_{C^{1,\alpha}}\le c\,\kappa_S\|h\|_{C^{3}}$ (L2(ii), $j\le2$,
plus interpolation to the $\alpha$-grade), so $a=2$,
$\|\Phi\|_{C^{1,\alpha}}\le c\,c_\theta\Lambda_\Delta\kappa_S
\|h\|_{C^3}\,|u(x)|\,\|v\|_{C^{1,\alpha}}$; apply (P2) for each
$\tau$ (the phase family $\{S_\tau\}$ is admissible by L2(iii)) and
integrate $d\tau$.
(C) $v^{(2)}_g-v^{(2)}_{g'}=\sum_{k\le2}\bigl[(v_k[g]-v_k[g'])
\sigma_g^k+v_k[g']\,(\sigma_g^k-\sigma_{g'}^k)\bigr]$ and
$\sigma_g^k-\sigma_{g'}^k=w\sum_{j<k}\sigma_g^j\sigma_{g'}^{k-1-j}$,
so $\|v^{(2)}_g-v^{(2)}_{g'}\|_{C^0(V)}\le
\kappa_C'\|h\|_{C^6}$ with
$\kappa_C':=c\bigl[\kappa_v(1+\Lambda_\sigma\delta_c^2)^2
+\Lambda_v\kappa_\sigma\delta_c^2(1+\Lambda_\sigma\delta_c^2)\bigr]$;
apply (P3), $a=0$, weight $|u(x)|\|v\|_{C^0}$.
(D) $F=v^{(2)}_{g'}w\,\theta_a u(x)v$: $a=2$,
$\|\Phi\|_{C^{0,\alpha}}\le c\,c_\theta\Lambda_v
(1+\Lambda_\sigma\delta_c^2)^2\kappa_S\|h\|_{C^2}\,|u(x)|
\,\|v\|_{C^{0,\alpha}}$; apply (P1) with $a=2$, per $\tau$.
Summing the four legs, the charts ($\le m(K)$ overlapping terms), and
integrating $|u(x)|$ over $K$:
\begin{equation}\label{port:eq-kappa}
  \bigl|(H_g-H_{g'})(u\otimes v)\bigr|
  \le\kappa(g_0,K)\,\|h\|_{C^6(K'')}\|u\|_{C^0(K)}\|v\|_{C^{1,\alpha}(K)},
\end{equation}
\[
  \kappa(g_0,K):=\frac{|K|\,m(K)\,c_\theta\,\kappa_P}{8\pi^2}
  \Bigl[\underbrace{\kappa_\Delta\tfrac{\delta_c^2}{2}}_{\mathrm{(A)}}
  +\underbrace{c\,\Lambda_\Delta\kappa_S\tfrac{\delta_c^2}{2}}
  _{\mathrm{(B)}}
  +\underbrace{\kappa_C'\,\delta_c^{2}
  \bigl(1+|\log\tfrac{\delta_c^2}{\ell^2}|\bigr)}_{\mathrm{(C)}}
  +\underbrace{c\,\Lambda_v(1+\Lambda_\sigma\delta_c^2)^2\kappa_S
  \tfrac{\delta_c^{4}}{4}}_{\mathrm{(D)}}\Bigr].
\]
This is the sharper (asymmetric) form. Since
$\|h\|_{C^6}\le\|h\|_{C^{6,\alpha}}\le C_{\mathrm{sob}}\|h\|_{H^s}$,
$\|u\|_{C^0}\le C_{\mathrm{sob}}\|u\|_{H^{s-2}}$ and
$\|v\|_{C^{1,\alpha}}\le C_{\mathrm{sob}}\|v\|_{H^{s-2}}$
\eqref{port:eq-csob}, the symmetric form \eqref{port:eq-clausei}
follows with $C_H=C_{\mathrm{sob}}^3\kappa$. All bounds are uniform in
$\epsilon\in(0,1]$; letting $\epsilon\downarrow0$ against the
per-metric limits (cited footing) completes the proof. No state was
used at any step. \qedhere
\end{proof}

\begin{remark}[Three-way fidelity: what this port states, against the
appendix demand and the established theorem]\label{port:rem-match}
(a) \emph{Appendix demand.} Display \eqref{port:eq-clausei} is verbatim
\eqref{eq:an17-clausei} $=$ clause (i) of Hyp.~\ref{hyp:hadamard-uniform} (Option B)
and of Hyp.~\ref{hyp:hadamard-uniform-omega}: same pairing, same three norms
$H^s,H^{s-2},H^{s-2}$, same support condition, and the quantifier is
over the \emph{pair} $(g,g')$ ranging over the full ball $U$ ---
strictly containing the $U_\omega\times U_\omega$ range item
(M3) asks for (Lem.~\ref{an17:lem-embed}). $H_g$ is the appendix's own object
\eqref{port:eq-Hdef} (the $N{=}2$ glued parametrix, one global $\ell$).
No hypothesis is strengthened; the only weakening is in the
\emph{conclusion's} favor (full ball, asymmetric sharper form).
(b) \emph{The established clause-(i) theorem.}
The constant reproduces its architecture exactly ---
$C_H=C_{\mathrm{sob}}^3\kappa$ with the cube accounting for the three
norm conversions \eqref{port:eq-csob}; the Christoffel difference is
the exact identity \eqref{port:eq-gammaid}; all transport legs are
Gronwall along geodesics (L2, L3); the $r^{-2}$ kernel enters as the
$a{=}2$, $p{=}2$ leg of L4, $L^1_{\mathrm{loc}}$ at $n=4$ with mass
$M_2$ \eqref{port:eq-masses}; the result is state-independent. The
suppressed arguments of the compact notation are displayed here:
$\kappa=\kappa(g_0,K;s,\varepsilon_0,\delta_c,\ell,\text{atlas})$.
(c) \emph{Single-metric consistency.} Applying L4 legs (P1)/(P3) to
$H_g$ alone (weights $\Lambda_\Delta,\Lambda_v$ instead of the
$\kappa$-differences) re-proves the single-metric
Lem.~\ref{lem:an17-clausei-single} bound with $C_H(K,g)\le
C_{\mathrm{sob}}^2\,\kappa'(g_0,K)$, consistent with GT1.
\end{remark}

\begin{remark}[Sharpening the test-space threshold: $t>3$, not
$t>2$]\label{port:rem-testspace}
A prior side note records that the $H^{s-2}$ test space is
over-strong and ``any $H^t$, $t>2$ suffices''. The re-derivation
\emph{confirms the over-strength but corrects the threshold}: the legs
(A), (C), (D) need only $C^{0,\alpha}$ ($t>2+\alpha$) in the $v$-slot
and $C^0$ in the $u$-slot, but the leading leg (B) --- the $r^{-2}$
kernel --- consumes one full transversal derivative of the test pair:
$\mathrm{Fp}(\sigma_\tau^{-2})$ pairs at the Lipschitz rate only
against $C^{1,\alpha}$, i.e.\ $H^t$, $t>3+\alpha$, in one slot. At
$2<t\le3$ the pairing rate degrades to H\"older
$\|g-g'\|^{t-2}$: a shifted first-order pole tested against
$C^{0,\beta}$ moves at rate $d^{\beta}$, not $d$ (numerically
confirmed at exponents $0.50$ vs $1.00$; the mechanism is the same
H\"older-vs-Lipschitz conversion loss as the EP2 fence, now on the
real side). Consequence for consumers: none --- the clause consumes
$H^{s-2}$, $s-2>6$, with room to spare on every leg; but the
side note should read ``any $H^t$, $t>3$, in one slot ($t>2$ in the
other)''. This corrects a side \emph{remark}, not any
printed statement of the paper or appendix.
\end{remark}

\begin{remark}[Grade and consumers; what this port does NOT do]
\label{port:rem-grade}
This port closes item (M3) of \S\ref{sec:an17-modulo} at THEOREM grade: the difference form
\eqref{port:eq-clausei} is now \emph{derived}, ball-uniform over the
family, at the hypothesis' own standing fence $s>n/2+6$ (no re-fence:
the proof consumes $\|h\|_{C^6}$, $\|g\|_{C^{6,\alpha}}$ with
$\alpha=\min(\tfrac12,\tfrac{s-8}2)$, both inside $H^{s}$, $s>8$).
Consumer effect (per the printed audit):
clause (i)-diff feeds only the singular leg of
Thm.~\ref{thm:E2-Lipschitz} and no quantitative slice-register consumer, so no
downstream constant changes; the pivot's conditionality on (M1) and
(M2) is untouched. This port says \emph{nothing} about the finite part
$W_g$ (clause (ii)/M1: still PLAUSIBLE at its printed modulo-list) or
the QEI constants (clause (iii)/M2: reduced-at-consumed-strength,
named). The paper's hypothesis statements are unchanged; this port supplies
the clause-(i) derivation only.
\end{remark}
% ====================== end of P-M3 fragment ==========================
